\documentclass[10pt,a4paper]{amsart}
\usepackage{amsmath,amssymb,amsthm,graphicx}
% Retained preamble for the Stacks Project chapter {CH} (shared project preamble).
% The chapter file pinned-original-source.tex only carries \input{{preamble}} on its first line,
% so no preamble text lives inside the chapter itself.  The Stacks Project uses one single shared
% preamble for every chapter: the snapshot stages it three times (sources/stacks-chow/preamble.tex,
% sources/stacks-etale-cohomology/preamble.tex and focused-seeds/AGP/sources/stacks-snapshot-20261011/
% preamble.tex) and those copies are identical in size.  The block below is copied mechanically,
% line for line, from sources/stacks-chow/preamble.tex lines 14-18 and 154-217: the Xy-pic setup and
% the native theorem declarations together with the mathematical macros.  Nothing was added, removed
% or redefined; the wrapper class and the amsmath/amssymb/amsthm/graphicx base packages are supplied
% by the assembler.
% For commutative diagrams we use Xy-pic
% Chapter: stacks-spaces-more-morphisms.tex -- More on Morphisms of Spaces.
\usepackage[all]{xy}

% We use 2cell for 2-commutative diagrams.
\xyoption{2cell}

% Theorem environments.
%
\theoremstyle{plain}
\newtheorem{theorem}[subsection]{Theorem}
\newtheorem{proposition}[subsection]{Proposition}
\newtheorem{lemma}[subsection]{Lemma}

\theoremstyle{definition}
\newtheorem{definition}[subsection]{Definition}
\newtheorem{example}[subsection]{Example}
\newtheorem{exercise}[subsection]{Exercise}
\newtheorem{situation}[subsection]{Situation}

\theoremstyle{remark}
\newtheorem{remark}[subsection]{Remark}
\newtheorem{remarks}[subsection]{Remarks}

\numberwithin{equation}{subsection}

% Macros
%
\def\lim{\mathop{\mathrm{lim}}\nolimits}
\def\colim{\mathop{\mathrm{colim}}\nolimits}
\def\Spec{\mathop{\mathrm{Spec}}}
\def\Hom{\mathop{\mathrm{Hom}}\nolimits}
\def\Ext{\mathop{\mathrm{Ext}}\nolimits}
\def\SheafHom{\mathop{\mathcal{H}\!\mathit{om}}\nolimits}
\def\SheafExt{\mathop{\mathcal{E}\!\mathit{xt}}\nolimits}
\def\Sch{\mathit{Sch}}
\def\Mor{\mathop{\mathrm{Mor}}\nolimits}
\def\Ob{\mathop{\mathrm{Ob}}\nolimits}
\def\Sh{\mathop{\mathit{Sh}}\nolimits}
\def\NL{\mathop{N\!L}\nolimits}
\def\CH{\mathop{\mathrm{CH}}\nolimits}
\def\proetale{{pro\text{-}\acute{e}tale}}
\def\etale{{\acute{e}tale}}
\def\QCoh{\mathit{QCoh}}
\def\Ker{\mathop{\mathrm{Ker}}}
\def\Im{\mathop{\mathrm{Im}}}
\def\Coker{\mathop{\mathrm{Coker}}}
\def\Coim{\mathop{\mathrm{Coim}}}

% Boxtimes
%
\DeclareMathSymbol{\boxtimes}{\mathbin}{AMSa}{"02}

%
% Macros for moduli stacks/spaces
%
\def\QCohstack{\mathcal{QC}\!\mathit{oh}}
\def\Cohstack{\mathcal{C}\!\mathit{oh}}
\def\Spacesstack{\mathcal{S}\!\mathit{paces}}
\def\Quotfunctor{\mathrm{Quot}}
\def\Hilbfunctor{\mathrm{Hilb}}
\def\Curvesstack{\mathcal{C}\!\mathit{urves}}
\def\Polarizedstack{\mathcal{P}\!\mathit{olarized}}
\def\Complexesstack{\mathcal{C}\!\mathit{omplexes}}
% \Pic is the operator that assigns to X its picard group, usage \Pic(X)
% \Picardstack_{X/B} denotes the Picard stack of X over B
% \Picardfunctor_{X/B} denotes the Picard functor of X over B
\def\Pic{\mathop{\mathrm{Pic}}\nolimits}
\def\Picardstack{\mathcal{P}\!\mathit{ic}}
\def\Picardfunctor{\mathrm{Pic}}
\def\Deformationcategory{\mathcal{D}\!\mathit{ef}}
\usepackage{cleveref}
\crefname{originalreference}{original reference}{original references}
\Crefname{originalreference}{Original reference}{Original references}
\makeatletter
\newlabel{coherent-remark-inverse-systems-kernel-cokernel-annihilated-by}{{0888}{source}{}{sourcepointer-1}{}}
\newlabel{coherent-remark-inverse-systems-kernel-cokernel-annihilated-by@cref}{{[originalreference][0][]0888}{[1][0][]source}{}{}{}}
\makeatother
\crefname{originalequation}{original equation}{original equations}
\Crefname{originalequation}{Original equation}{Original equations}
\makeatletter
\newlabel{equation-completion-functor-proper-over-A}{{08B8}{source}{}{sourcepointer-2}{}}
\newlabel{equation-completion-functor-proper-over-A@cref}{{[originalequation][0][]08B8}{[1][0][]source}{}{}{}}
\makeatother
\crefname{originallemma}{original lemma}{original lemmas}
\Crefname{originallemma}{Original lemma}{Original lemmas}
\makeatletter
\newlabel{lemma-existence-easy}{{08BB}{source}{}{sourcepointer-3}{}}
\newlabel{lemma-existence-easy@cref}{{[originallemma][0][]08BB}{[1][0][]source}{}{}{}}
\makeatother
\crefname{originalreference}{original reference}{original references}
\Crefname{originalreference}{Original reference}{Original references}
\makeatletter
\newlabel{situation-existence}{{08B7}{source}{}{sourcepointer-4}{}}
\newlabel{situation-existence@cref}{{[originalreference][0][]08B7}{[1][0][]source}{}{}{}}
\makeatother
\crefname{originalreference}{original reference}{original references}
\Crefname{originalreference}{Original reference}{Original references}
\makeatletter
\newlabel{spaces-cohomology-lemma-i-star-equivalence}{{08AM}{source}{}{sourcepointer-5}{}}
\newlabel{spaces-cohomology-lemma-i-star-equivalence@cref}{{[originalreference][0][]08AM}{[1][0][]source}{}{}{}}
\makeatother
\crefname{originalreference}{original reference}{original references}
\Crefname{originalreference}{Original reference}{Original references}
\makeatletter
\newlabel{spaces-perfect-lemma-union-closed-proper-over-base}{{0CZH}{source}{}{sourcepointer-6}{}}
\newlabel{spaces-perfect-lemma-union-closed-proper-over-base@cref}{{[originalreference][0][]0CZH}{[1][0][]source}{}{}{}}
\makeatother
\title{A refined algebraization step: gluing coherent systems with proper support}
\author{The Stacks Project Authors}
\date{stacks-project:spaces-more-morphisms.tex:local-snapshot-20261011; GNU FDL 1.2 or later}
\begin{document}
\maketitle
\noindent\textit{Editorial note.} Counted span is the author's own proof environment, source lines 9849-9938, immediately following the statement at 9827-9847.\par\medskip
\noindent\textit{Editorial note.} The statement itself references the situation of the section by \textbackslash{}ref, and the proof cites several cross-chapter results; all are navigation pointers recorded in assembly.external\_references.\par\medskip
\noindent\textit{Editorial note.} The final properness argument is displayed in full inside the span (kernel support, G\_1 support, H support, union lemma).\par\medskip
\noindent\textit{Editorial note.} Steps the author leaves to the reader are flagged here for the reviewer: source line(s) 9897, 9906 carry an "omitted"/"follows easily"/"immediate from the definitions" phrase. Those sentences are retained verbatim; the displayed argument around them stays inside the counted span.\par\medskip
\noindent\textit{Allowed standard prerequisites (editorial).}\begin{itemize}
\item The situation of Grothendieck's existence theorem. chapter situation situation-existence, cited in the statement and by equation-completion-functor-proper-over-A. Fixes the formal setup: the ideal I, the thickenings X\_e, and the category Coh with support proper over A.
\item Easy algebraization step. chapter lemma lemma-existence-easy, cited inside the proof. Provides the simpler algebraization statement that the refined argument reduces to.
\item Inverse systems with kernel and cokernel annihilated by I. Coherent, remark coherent-remark-inverse-systems-kernel-cokernel-annihilated-by, cited inside the proof. Controls how a map of inverse systems behaves on kernels and cokernels under the thickenings.
\item Union of closed subsets proper over the base. Derived Categories of Spaces, lemma spaces-perfect-lemma-union-closed-proper-over-base, cited at proof lines 9936-9937. Concludes properness of the support of the algebraized module from properness of the supports of its kernel and of H.
\end{itemize}\medskip
\ifdefined\hypertarget\hypertarget{sourcepointer-1}{}\fi
\noindent\textit{Source reference.} Original reference 0888: focused-seeds/geometry/S9-spaces-existence-tricky/ref-tags.txt, source lines 4--4. Navigation pointer only: the counted proof cites this statement by \textbackslash{}ref as a standing input; the deduction the counted proof has to make is carried out inside the counted span, and no part of that deduction is delegated to this pointer.\par\medskip
\ifdefined\hypertarget\hypertarget{sourcepointer-2}{}\fi
\noindent\textit{Source reference.} Original equation 08B8: sources/stacks-spaces-more-morphisms/pinned-original-source.tex, source lines 9710--9710. Navigation pointer only: the counted proof cites this statement by \textbackslash{}ref as a standing input; the deduction the counted proof has to make is carried out inside the counted span, and no part of that deduction is delegated to this pointer.\par\medskip
\ifdefined\hypertarget\hypertarget{sourcepointer-3}{}\fi
\noindent\textit{Source reference.} Original lemma 08BB: sources/stacks-spaces-more-morphisms/pinned-original-source.tex, source lines 9801--9801. Navigation pointer only: the counted proof cites this statement by \textbackslash{}ref as a standing input; the deduction the counted proof has to make is carried out inside the counted span, and no part of that deduction is delegated to this pointer.\par\medskip
\ifdefined\hypertarget\hypertarget{sourcepointer-4}{}\fi
\noindent\textit{Source reference.} Original reference 08B7: sources/stacks-spaces-more-morphisms/pinned-original-source.tex, source lines 9682--9682. Navigation pointer only: the counted proof cites this statement by \textbackslash{}ref as a standing input; the deduction the counted proof has to make is carried out inside the counted span, and no part of that deduction is delegated to this pointer.\par\medskip
\ifdefined\hypertarget\hypertarget{sourcepointer-5}{}\fi
\noindent\textit{Source reference.} Original reference 08AM: focused-seeds/geometry/S9-spaces-existence-tricky/ref-tags.txt, source lines 5--5. Navigation pointer only: the counted proof cites this statement by \textbackslash{}ref as a standing input; the deduction the counted proof has to make is carried out inside the counted span, and no part of that deduction is delegated to this pointer.\par\medskip
\ifdefined\hypertarget\hypertarget{sourcepointer-6}{}\fi
\noindent\textit{Source reference.} Original reference 0CZH: focused-seeds/geometry/S9-spaces-existence-tricky/ref-tags.txt, source lines 6--6. Navigation pointer only: the counted proof cites this statement by \textbackslash{}ref as a standing input; the deduction the counted proof has to make is carried out inside the counted span, and no part of that deduction is delegated to this pointer.\par\medskip
\begin{lemma}
\label{lemma-existence-tricky}
In Situation \ref{situation-existence}. Let $\mathcal{K} \subset \mathcal{O}_X$
be a quasi-coherent sheaf of ideals. Let $X_e \subset X$ be the closed subspace
cut out by $\mathcal{K}^e$. Let $\mathcal{I}_e = \mathcal{I}\mathcal{O}_{X_e}$.
Let $(\mathcal{F}_n)$ be an object of
$\textit{Coh}_{\text{support proper over } A}(X, \mathcal{I})$.
Assume
\begin{enumerate}
\item the functor
$\textit{Coh}_{\text{support proper over } A}(\mathcal{O}_{X_e})
\to \textit{Coh}_{\text{support proper over } A}(X_e, \mathcal{I}_e)$
is an equivalence for all $e \geq 1$, and
\item there exists an object $\mathcal{H}$ of
$\textit{Coh}_{\text{support proper over } A}(\mathcal{O}_X)$ and a map
$\alpha : (\mathcal{F}_n) \to \mathcal{H}^\wedge$ whose
kernel and cokernel are annihilated by a power of $\mathcal{K}$.
\end{enumerate}
Then $(\mathcal{F}_n)$ is in the essential image of
(\ref{equation-completion-functor-proper-over-A}).
\end{lemma}

\begin{proof}
During this proof we will use without further mention that for a closed
immersion $i : Z \to X$ the functor $i_*$ gives an equivalence between the
category of coherent modules on $Z$ and coherent modules on $X$ annihilated
by the ideal sheaf of $Z$, see
Cohomology of Spaces, Lemma \ref{spaces-cohomology-lemma-i-star-equivalence}.
In particular we think of
$$
\textit{Coh}_{\text{support proper over } A}(\mathcal{O}_{X_e})
\subset
\textit{Coh}_{\text{support proper over } A}(\mathcal{O}_X)
$$
as the full subcategory of consisting of modules annihilated by
$\mathcal{K}^e$ and
$$
\textit{Coh}_{\text{support proper over } A}(X_e, \mathcal{I}_e)
\subset
\textit{Coh}_{\text{support proper over } A}(X, \mathcal{I})
$$
as the full subcategory of objects annihilated by $\mathcal{K}^e$.
Moreover (1) tells us these two categories are equivalent under the
completion functor (\ref{equation-completion-functor-proper-over-A}).

\medskip\noindent
Applying this equivalence we get a coherent $\mathcal{O}_X$-module
$\mathcal{G}_e$ annihilated by $\mathcal{K}^e$ corresponding to the system
$(\mathcal{F}_n/\mathcal{K}^e\mathcal{F}_n)$ of
$\textit{Coh}_{\text{support proper over } A}(X, \mathcal{I})$. The maps
$\mathcal{F}_n/\mathcal{K}^{e + 1}\mathcal{F}_n \to
\mathcal{F}_n/\mathcal{K}^e\mathcal{F}_n$ correspond to canonical maps
$\mathcal{G}_{e + 1} \to \mathcal{G}_e$ which induce isomorphisms
$\mathcal{G}_{e + 1}/\mathcal{K}^e\mathcal{G}_{e + 1} \to \mathcal{G}_e$.
We obtain an object $(\mathcal{G}_e)$ of the category
$\textit{Coh}_{\text{support proper over } A}(X, \mathcal{K})$.
The map $\alpha$ induces a system of maps
$$
\mathcal{F}_n/\mathcal{K}^e\mathcal{F}_n
\longrightarrow
\mathcal{H}/(\mathcal{I}^n + \mathcal{K}^e)\mathcal{H}
$$
whence maps $\mathcal{G}_e \to \mathcal{H}/\mathcal{K}^e\mathcal{H}$
(by the equivalence of categories again).
Let $t \geq 1$ be an integer, which exists by assumption (2),
such that $\mathcal{K}^t$ annihilates the kernel and cokernel of all the maps
$\mathcal{F}_n \to \mathcal{H}/\mathcal{I}^n\mathcal{H}$.
Then $\mathcal{K}^{2t}$ annihilates the kernel and cokernel of the maps
$\mathcal{F}_n/\mathcal{K}^e\mathcal{F}_n \to
\mathcal{H}/(\mathcal{I}^n + \mathcal{K}^e)\mathcal{H}$
(details omitted; see Cohomology of Schemes,
Remark \ref{coherent-remark-inverse-systems-kernel-cokernel-annihilated-by}).
Whereupon we conclude that $\mathcal{K}^{4t}$ annihilates the kernel and
the cokernel of the maps
$$
\mathcal{G}_e
\longrightarrow
\mathcal{H}/\mathcal{K}^e\mathcal{H},
$$
(details omitted;  see Cohomology of Schemes,
Remark \ref{coherent-remark-inverse-systems-kernel-cokernel-annihilated-by}).
We apply Lemma \ref{lemma-existence-easy} to obtain a coherent
$\mathcal{O}_X$-module $\mathcal{F}$, a map
$a : \mathcal{F} \to \mathcal{H}$ and an isomorphism
$\beta : (\mathcal{G}_e) \to (\mathcal{F}/\mathcal{K}^e\mathcal{F})$
in $\textit{Coh}(X, \mathcal{K})$. Working backwards, for a given $n$
the triple
$(\mathcal{F}/\mathcal{I}^n\mathcal{F}, a \bmod \mathcal{I}^n, \beta
\bmod \mathcal{I}^n)$ is a triple as in the lemma for the morphism
$\alpha_n \bmod \mathcal{K}^e :
(\mathcal{F}_n/\mathcal{K}^e\mathcal{F}_n) \to
(\mathcal{H}/(\mathcal{I}^n + \mathcal{K}^e)\mathcal{H})$
of $\textit{Coh}(X, \mathcal{K})$. Thus the uniqueness in
Lemma \ref{lemma-existence-easy}
gives a canonical isomorphism
$\mathcal{F}/\mathcal{I}^n\mathcal{F} \to \mathcal{F}_n$
compatible with all the morphisms in sight.

\medskip\noindent
To finish the proof of the lemma we still have to show that the
support of $\mathcal{F}$ is proper over $A$.
By construction the kernel of $a : \mathcal{F} \to \mathcal{H}$
is annihilated by a power of $\mathcal{K}$. Hence the support of
this kernel is contained in the support of $\mathcal{G}_1$. Since
$\mathcal{G}_1$ is an object of
$\textit{Coh}_{\text{support proper over } A}(\mathcal{O}_{X_1})$
we see this is proper over $A$. Combined with the fact that the
support of $\mathcal{H}$ is proper over $A$ we conclude that the
support of $\mathcal{F}$ is proper over $A$ by
Derived Categories of Spaces, Lemma
\ref{spaces-perfect-lemma-union-closed-proper-over-base}.
\end{proof}

\end{document}
